\documentclass[conference]{IEEEtran}
\IEEEoverridecommandlockouts

\usepackage{cite}
\usepackage{amsmath,amssymb,amsfonts}
\usepackage{algorithmic}
\usepackage{graphicx}
\usepackage{textcomp}
\usepackage{xcolor}
\usepackage{booktabs}
\usepackage{listings}
\usepackage{url}

% Code listing style
\lstset{
  basicstyle=\ttfamily\footnotesize,
  keywordstyle=\bfseries,
  commentstyle=\itshape\color{gray},
  breaklines=true,
  frame=single,
  xleftmargin=6pt,
  xrightmargin=6pt,
  framesep=4pt,
  language=C++
}

\def\BibTeX{{\rm B\kern-.05em{\sc i\kern-.025em b}\kern-.08em
    T\kern-.1667em\lower.7ex\hbox{E}\kern-.125emX}}

\begin{document}

\title{GPU Modular Arithmetic Unit with Barrett Reduction}

\author{%
\IEEEauthorblockN{Akhil Teja Chirra}
\IEEEauthorblockA{\textit{Dept. of Electrical and}\\
\textit{Computer Engineering}\\
UC San Diego\\
achirra@ucsd.edu}
\and
\IEEEauthorblockN{Vagda Sameera Talari}
\IEEEauthorblockA{\textit{Dept. of Electrical and}\\
\textit{Computer Engineering}\\
UC San Diego\\
vtalari@ucsd.edu}
\and
\IEEEauthorblockN{Abhijit Gupta}
\IEEEauthorblockA{\textit{Dept. of Electrical and}\\
\textit{Computer Engineering}\\
UC San Diego\\
akg008@ucsd.edu}
}

\maketitle

\begin{abstract}
We implement a prime-agnostic modular arithmetic engine over large prime fields
entirely from scratch in pure CUDA, targeting an NVIDIA Tesla T4 GPU (sm\_75).
The engine supports modular multiplication, modular exponentiation, and Barrett
reduction, applied to a Number Theoretic Transform (NTT)---the modular-arithmetic
core of lattice-based post-quantum cryptography (Kyber, Dilithium) and
zero-knowledge proof systems. We benchmark GPU implementations against
single-threaded CPU equivalents using the identical Barrett C++ code path.
Measured results on a Colab T4: Barrett modmul achieves a \textbf{569$\times$
GPU speedup} (131~Gops/s), the NTT kernel peaks at \textbf{52$\times$} speedup
at $n=2^{18}$, and batch modular exponentiation over the 64-bit Goldilocks prime
achieves \textbf{277$\times$} speedup with zero correctness failures. We analyze
the performance-security trade-offs including timing side-channel exposure and
the compute-to-memory-bandwidth transition at scale. Code is available at
\url{https://github.com/achirra-cipher/ece268-gpu-modarith}.
\end{abstract}

\begin{IEEEkeywords}
Barrett reduction, modular arithmetic, GPU cryptography, CUDA, Number Theoretic
Transform, post-quantum cryptography, NTT, modular exponentiation
\end{IEEEkeywords}

%%--------------------------------------------------------------------
\section{Introduction}
%%--------------------------------------------------------------------

Modern cryptographic workloads---lattice-based post-quantum schemes
(Kyber~\cite{kyber}, Dilithium~\cite{dilithium}) and zero-knowledge proof
systems (Plonky2~\cite{plonky2})---are dominated by modular arithmetic over
large prime fields. The Number Theoretic Transform (NTT), the algebraic analogue
of the FFT over a finite field, accounts for the majority of compute in these
schemes. An NTT of size $n$ consists of $\log_2 n$ stages, each with $n/2$
independent butterfly operations; the workload is massively data-parallel and
each butterfly requires two modular multiplications, making NTT throughput
directly proportional to modmul throughput.

GPUs expose thousands of parallel cores and are natural accelerators for such
workloads. Prior work has demonstrated GPU acceleration of elliptic-curve scalar
multiplication, hash functions, and symmetric ciphers, but modular arithmetic
engines targeting the prime fields used in post-quantum lattice schemes and
ZK-proof systems remain less explored in the pure-CUDA setting. A key challenge
is modular reduction: standard integer division is slow on GPUs---the hardware
lacks a native 64-bit integer divide instruction---and data-dependent division
latency introduces a potential timing side-channel in cryptographic contexts.
Barrett reduction~\cite{barrett} addresses both problems by replacing division
with a multiply-shift sequence using a precomputed constant~$\mu$, making the
reduction branch-light, divide-free, and GPU-efficient.

This paper presents a pure-CUDA modular arithmetic engine that implements:
(i)~Barrett reduction for two prime fields,
(ii)~modular multiplication (\texttt{modmul}),
(iii)~modular exponentiation (\texttt{modexp}) via square-and-multiply, and
(iv)~modular inverse (\texttt{modinv}) via Fermat's little theorem.
The engine is exercised through a GPU NTT and a batch modexp demo over the
64-bit Goldilocks prime. We benchmark GPU versus single-threaded CPU performance
and analyze correctness, throughput, and security trade-offs.
The entire implementation is written in pure CUDA/C++ with no PyCUDA, CuPy, or
other Python GPU wrappers, qualifying for the bonus pure-CUDA track.

%%--------------------------------------------------------------------
\section{Background}
%%--------------------------------------------------------------------

\subsection{Barrett Reduction}

Given an integer $x \in [0,\,p^2)$ and a modulus $p$, Barrett reduction computes
$x \bmod p$ without division. A constant $\mu = \lfloor 2^k / p \rfloor$ is
precomputed offline. The reduction then proceeds as:
\begin{align}
  q &= (x \cdot \mu) \gg k  \label{eq:q} \\
  r &= x - q \cdot p         \label{eq:r} \\
  r &\mathrel{-}= p \quad \text{if } r \geq p \label{eq:corr}
\end{align}
Step~\eqref{eq:q} estimates $\lfloor x/p \rfloor$, which is off by at most~1.
Step~\eqref{eq:r} produces $r \in [0,\,2p)$; the single conditional subtraction
in~\eqref{eq:corr} brings $r$ into $[0,\,p)$. This reduces modular reduction to
two multiplications and one shift---no hardware division required.

The constant $\mu$ is shared across all threads and resides in constant cache or
a register broadcast. On sm\_75 (T4), the 64-bit high-half multiply maps to the
\texttt{\_\_umul64hi} hardware intrinsic---a single instruction---making the hot
path two 64-bit multiplies, cheap and warp-uniform with no thread divergence from
the reduction itself.

\subsection{Number Theoretic Transform}

The NTT over $\mathbb{F}_p$ requires $p \equiv 1 \pmod{n}$ so that $n$-th roots
of unity exist in the field~\cite{nttref}. The Cooley-Tukey butterfly at each stage computes:
\begin{align}
  A[u] &= (u + \omega \cdot v) \bmod p \label{eq:but1} \\
  A[v] &= (u - \omega \cdot v) \bmod p \label{eq:but2}
\end{align}
where $\omega$ is a twiddle factor. Each butterfly requires two modular
multiplications (\eqref{eq:but1}--\eqref{eq:but2}), so NTT throughput is
directly proportional to modmul throughput. The $\log_2 n$ stages each consist
of $n/2$ independent butterfly pairs, making the NTT ideal for GPU parallelism.

\subsection{Barrett vs.\ Montgomery}

Montgomery multiplication~\cite{montgomery} is an alternative that amortizes
the division cost over multiple multiplications by working in a Montgomery domain.
It typically requires one fewer conditional subtraction per reduction in tight
modexp loops, because the correction term is naturally absorbed into the
multiply-and-add chain. However, Montgomery requires explicit conversion of
operands into Montgomery form ($a \to aR \bmod p$, where $R = 2^{64}$) and
back out afterward. For standalone modular exponentiation this conversion cost
amortizes over many multiplications on the same value.

For NTT, Barrett is the better choice: the $n$ input elements must stay in
natural residue form across all $\log_2 n$ butterfly stages. Switching to
Montgomery would require two additional O($n$) passes---one to convert the
entire array into Montgomery domain before stage 1, and one to convert back
after the final stage. As we show in Section~VI, the NTT is already
memory-bandwidth-bound at large $n$; these two extra passes would directly reduce
effective speedup. Our engine is prime-agnostic: any prime $p$ requires only
recomputing $\mu = \lfloor 2^k / p \rfloor$ with no other structural changes.

%%--------------------------------------------------------------------
\section{Implementation}
%%--------------------------------------------------------------------

\subsection{Prime Field Selection}

We use two prime fields (Table~\ref{tab:primes}).
\textbf{P32} (998{,}244{,}353 $= 119 \cdot 2^{23} + 1$) is the primary NTT
benchmark prime: 30-bit, NTT-friendly (2-adicity~23, generator $g=3$), with
all products fitting in 64 bits, enabling Barrett with native
\texttt{\_\_uint128\_t}. \textbf{P64} ($2^{64} - 2^{32} + 1$, the Goldilocks
prime) is used in real zero-knowledge proof systems (Plonky2) and requires
128-bit intermediates on device---the bonus large-prime path.

\begin{table}[t]
\caption{Prime Fields Used}
\label{tab:primes}
\begin{center}
\begin{tabular}{|c|l|c|l|}
\hline
\textbf{Name} & \textbf{Prime} & \textbf{Bits} & \textbf{Role} \\
\hline
P32 & $119 \cdot 2^{23}+1$ & 30 & Primary (NTT) \\
\hline
P64 & $2^{64}-2^{32}+1$ & 64 & Bonus (modexp, zk) \\
\hline
\end{tabular}
\end{center}
\end{table}

\subsection{Core Arithmetic: \texttt{modarith.cuh}}

All arithmetic lives in a single CUDA header \texttt{modarith.cuh} marked
\texttt{\_\_host\_\_ \_\_device\_\_}, enabling the same Barrett and modmul
logic to compile and run both on the CPU (for local correctness testing) and
on the GPU (for production). The key device-side optimization uses the
\texttt{\_\_umul64hi} intrinsic---a single sm\_75 instruction that computes
the high 64 bits of a 128-bit product:

\begin{lstlisting}
// Pseudocode — illustrates algorithm structure
// Actual implementation: cuda/include/modarith.cuh

__host__ __device__ static inline u64
barrett32_reduce(u64 x, const Barrett32& B) {
  // q = floor(x / p), off by at most 1
  u64 q = ((__uint128_t)x * B.mu) >> B.k;
  u64 r = x - q * B.p;           // r in [0,2p)
  if (r >= B.p) r -= B.p;        // single correction
  return r;
}

__device__ static inline u64
mulmod32(u64 a, u64 b, const Barrett32& B) {
  u64 hi = __umul64hi(a, b); // 1 sm_75 instruction
  u64 lo = a * b;
  return barrett32_reduce((hi << 32) | (lo >> 32), B);
}
\end{lstlisting}

Modular exponentiation is implemented via square-and-multiply, iterating
over the bits of the exponent (up to 64 bits for 64-bit exponents; 30 bits for
P32 with $p-2$ as exponent). Modular inverse uses Fermat's
little theorem: $\text{inv}(a) = a^{p-2} \bmod p$, implemented by calling
\texttt{modexp} with exponent $p-2$. Both operations are fully portable across
P32 and P64 via the same Barrett interface.

\subsection{GPU NTT Kernel}

The NTT is implemented as an iterative Cooley-Tukey transform, with one CUDA
thread handling one butterfly pair per stage:

\begin{lstlisting}
// Pseudocode — actual kernel: cuda/include/ntt_gpu.cuh
// Driver in cuda/src/ntt.cu

__global__ void k_ntt_stage(
    u64* A, int half, const u64* roots,
    int n, const Barrett32 B) {
  int tid = blockIdx.x * blockDim.x + threadIdx.x;
  if (tid >= n/2) return;
  int j   = tid % half;
  int u_i = (tid/half)*2*half + j;
  int v_i = u_i + half;
  u64 tw  = roots[(n / (2*half)) * j]; // precomputed
  u64 u   = A[u_i];
  u64 v   = mulmod32(A[v_i], tw, B);
  A[u_i]  = modadd(u, v, B.p);
  A[v_i]  = modsub(u, v, B.p);
}
\end{lstlisting}

The kernel is launched $\log_2 n$ times---once per NTT stage---with
$\lceil n/2 / 256 \rceil$ blocks of 256 threads each (block size chosen to
maximise occupancy on sm\_75 with its 2048 resident threads per SM). Each
launch implicitly synchronises the entire device, providing stage-to-stage
coherence without explicit barriers.

\textbf{Bit-reversal permutation.} Before the forward NTT, a separate kernel
rearranges the input array into bit-reversed order. Each thread $i$ swaps
$A[i]$ with $A[\text{bitrev}(i)]$ when $\text{bitrev}(i) > i$, computed with a
standard $\log_2 n$-step bit-flip loop. The permutation runs in a single kernel
launch and is counted in the ``GPU total'' timing column of Table~\ref{tab:ntt}.

\textbf{Inverse NTT (INTT).} The INTT reuses the same
\texttt{k\_ntt\_stage} kernel with inverse twiddle factors
$\omega^{-1} = \texttt{modpow}(\omega,\, p-1-n/p,\, B)$ and a final
element-wise scaling by $n^{-1} \bmod p$ (one modmul per element). This
design avoids code duplication: the forward and inverse transforms share the
identical butterfly implementation. Correctness is verified by the round-trip
identity $\text{INTT}(\text{NTT}(x)) = x$ for all tested sizes up to
$n = 2^{14}$ (Section~\ref{sec:val}).

Choosing Barrett over Montgomery avoids domain-conversion passes
between stages, which is critical because NTT becomes memory-bandwidth-bound at
large $n$.

\subsection{Repository Structure}

The complete codebase is organized as follows:
\begin{itemize}
  \item \texttt{cuda/include/modarith.cuh} --- Barrett, modmul, modexp, modinv (host+device)
  \item \texttt{cuda/include/params.h} --- Prime parameters and precomputed $\mu$
  \item \texttt{cuda/src/ntt.cu} --- GPU NTT butterfly kernels + bit-reversal
  \item \texttt{cuda/src/bench.cu} --- CUDA-event benchmark harness
  \item \texttt{cuda/src/modexp\_demo.cu} --- Batch modexp over P64
  \item \texttt{python/oracle.py} --- Ground-truth test vectors via sympy
  \item \texttt{colab/run\_on\_colab.ipynb} --- End-to-end Colab T4 notebook
\end{itemize}

\subsection{Development Workflow}

Development was conducted on an Apple Silicon Mac (arm64, no CUDA). The
\texttt{\_\_host\_\_ \_\_device\_\_} dual-compilation strategy allows the
identical Barrett and modmul code to be compiled as plain C++ with the host compiler (\texttt{clang} on Mac, \texttt{g++} on Linux/Colab)
and unit-tested locally against Python/sympy ground truth, with no GPU required.
Only after the CPU unit tests pass does the code move to GPU compilation.

The workflow proceeds in three phases. \emph{Phase 1 (local):} \texttt{host\_test.cpp}
includes \texttt{modarith.cuh} directly, runs 2M+ random test cases against
128-bit ground truth, and compares against \texttt{oracle.py} output. \emph{Phase 2
(Colab compile):} the \texttt{.cu} sources are compiled with \texttt{nvcc -arch=sm\_75
-O3} on a Colab T4 instance; if Phase 1 passed, compilation errors almost always
indicate a GPU-specific intrinsic issue rather than an arithmetic bug. \emph{Phase 3
(GPU verify):} the compiled binary runs the same test suite on GPU, and outputs
are compared byte-for-byte with the Phase 1 CPU reference. This three-phase
approach reduced GPU debugging time significantly, as the expensive Colab session
was only needed for performance measurement, not arithmetic validation.

%%--------------------------------------------------------------------
\section{Correctness \& Validation}\label{sec:val}
%%--------------------------------------------------------------------

We enforce a three-layer validation protocol:
GPU~C++ (nvcc) $\equiv$ CPU~C++ (clang) $\equiv$ Python (sympy). The protocol
is designed so that no two of the three layers share implementation code or
arithmetic libraries: the GPU and CPU paths share \texttt{modarith.cuh} but
differ in compiler and hardware; the Python layer uses sympy's arbitrary-precision
integer arithmetic, which is completely independent. Agreement across all three
layers therefore provides strong evidence of correctness.

\textbf{Barrett32 / Barrett64.} We generate 2 million+ random operand pairs
$(a, b)$ uniformly distributed in $[0, p)$ and verify that \texttt{modmul}$(a,
b)$ matches $a \cdot b \bmod p$ computed via \texttt{\_\_uint128\_t} ground truth.
Boundary operands ($a = 0$, $a = p-1$, products near $p^2$) are tested
separately. Both Barrett32 (P32) and Barrett64 (P64) pass with zero failures.

\textbf{modexp / modinv.} We test 100K random exponents, verifying
\texttt{modexp}$(a, e, p) = a^e \bmod p$ against an independent square-and-multiply
implemented with native 128-bit arithmetic. \texttt{modinv}$(a, p)$ is verified
via the identity $a \cdot \text{inv}(a) \equiv 1 \pmod{p}$.

\textbf{NTT correctness.} The CPU NTT is validated by the round-trip identity
INTT(NTT($x$)) $= x$ for $n$ up to $2^{14}$, and by verifying that
NTT-based cyclic convolution matches the naive O($n^2$) pointwise product for
$n$ up to $2^9$. The GPU forward NTT is then verified to match the CPU output
element-wise for all sizes $2^4$ through $2^{20}$---\emph{every element} of
every transform matches exactly.

\textbf{Batch modexp (GPU, P64).} We run 1,048,576 independent 64-bit
exponentiations over the Goldilocks prime on the T4 and compare element-wise
against the CPU reference. Zero mismatches across the full 1M-element batch.

%%--------------------------------------------------------------------
\section{Performance Results}
%%--------------------------------------------------------------------

\subsection{Experimental Setup}

All benchmarks run on a Colab NVIDIA Tesla T4 (sm\_75, 16~GB GDDR6,
2560~CUDA cores) versus a single CPU thread (Colab's Xeon, $\approx$2~GHz
effective clock). GPU timing uses \texttt{cudaEventRecord} /
\texttt{cudaEventElapsedTime} and excludes Python/host overhead. The CPU uses
the same Barrett C++ code compiled with the host C++ compiler (clang on Mac, \texttt{g++} on Colab).

\subsection{Barrett Modmul Throughput}

Table~\ref{tab:modmul} shows throughput for 4 million independent modular
multiplications with 64 iterations per lane, fully saturating all T4 CUDA
cores.

\begin{table}[t]
\caption{Barrett Modmul Throughput ($n=4$M, P32)}
\label{tab:modmul}
\begin{center}
\begin{tabular}{|l|c|c|c|}
\hline
\textbf{Metric} & \textbf{CPU (1 thread)} & \textbf{GPU kernel (T4)} & \textbf{Speedup} \\
\hline
Throughput & 0.23 Gops/s & 131.4 Gops/s & \textbf{569$\times$} \\
Kernel wall time & 1161 ms & 2.04 ms & 569$\times$ \\
\hline
\end{tabular}
\end{center}
\end{table}

The 569$\times$ speedup reflects the embarrassingly parallel nature of
independent modmul: one thread per element, no inter-thread communication, and
all 2560 T4 CUDA cores kept fully saturated.

\subsection{NTT Latency vs.\ Array Size}

Table~\ref{tab:ntt} shows NTT kernel latency and speedup across array sizes.

\begin{table}[t]
\caption{NTT Kernel Speedup (Forward NTT, P32)}
\label{tab:ntt}
\begin{center}
\begin{tabular}{|c|c|c|c|c|}
\hline
\textbf{$n$} & \textbf{GPU kernel} & \textbf{GPU total} & \textbf{CPU} & \textbf{Speedup} \\
\hline
$2^{10}$ & 0.051 ms & 0.100 ms & 0.046 ms & 0.9$\times$ \\
$2^{12}$ & 0.061 ms & 0.096 ms & 0.193 ms & 3.2$\times$ \\
$2^{14}$ & 0.079 ms & 0.163 ms & 0.866 ms & 11$\times$  \\
$2^{16}$ & 0.141 ms & 0.423 ms & 3.82 ms  & 27$\times$  \\
$2^{18}$ & 0.367 ms & 1.312 ms & 19.0 ms  & \textbf{52$\times$} \\
$2^{20}$ & 2.54 ms  & 6.37 ms  & 90.1 ms  & 35$\times$  \\
$2^{22}$ & 12.7 ms  & 27.3 ms  & 507 ms   & 40$\times$  \\
\hline
\end{tabular}
\end{center}
\end{table}

Three important observations emerge.
(i)~At $n=2^{10}$, the GPU kernel is \emph{slower} than a single CPU core
(0.9$\times$). The primary cause is kernel-launch latency: CUDA's
\texttt{cudaEventRecord}/\texttt{cudaEventElapsedTime} measurements reveal a
fixed per-launch overhead of roughly 40--60~$\mu$s on the Colab T4, which
dominates when useful work is only $n/2 = 512$ butterfly operations. The GPU
only overtakes the CPU around $n=2^{11}$--$2^{12}$, where $\log_2 n \approx 11$
stage launches together consume enough compute to amortize that fixed cost.
(ii)~Peak speedup (52$\times$, kernel-only) occurs at $n=2^{18}$ (256K
elements), the sweet spot where the working set fits in T4's L2 cache and all
20~SMs are saturated.
(iii)~Beyond $n=2^{18}$, speedup drops to $\approx35$--$40\times$ as the
working set overflows L2 and the kernel becomes memory-bandwidth-bound.

\subsection{Batch Modular Exponentiation (P64)}

Table~\ref{tab:modexp} shows results for 1M independent 64-bit modular
exponentiations over the Goldilocks prime.

\begin{table}[t]
\caption{Batch Modexp over P64 Goldilocks ($n = 2^{20}$)}
\label{tab:modexp}
\begin{center}
\begin{tabular}{|c|c|c|c|}
\hline
\textbf{CPU} & \textbf{GPU (T4)} & \textbf{Speedup} & \textbf{Mismatches} \\
\hline
955.8 ms & 3.4 ms & \textbf{277$\times$} & 0 \\
\hline
\end{tabular}
\end{center}
\end{table}

The 277$\times$ speedup is compute-bound: each thread independently runs 64
square-and-multiply iterations with no inter-thread communication and
near-uniform control flow---the archetypal GPU cryptographic workload.

%%--------------------------------------------------------------------
\section{Analysis}
%%--------------------------------------------------------------------

\subsection{Why Modmul Speedup ($569\times$) $\gg$ NTT Speedup ($52\times$)}\label{sec:bw}

The modmul benchmark is perfectly embarrassingly parallel---zero dependencies
between threads, perfectly coalesced memory access. The NTT butterfly pattern
introduces two sources of inefficiency: (i)~\emph{inter-stage dependencies}:
each of the $\log_2 n$ stages requires a separate kernel launch and a full pass
over device memory; (ii)~\emph{strided, non-coalesced global memory reads} in
early stages, where butterfly pairs are separated by large strides. At large $n$, the NTT transitions from compute-bound to
memory-bandwidth-bound, capping speedup well below the modmul case. The T4's
peak global memory bandwidth is approximately 320~GB/s; at $n=2^{22}$ each
stage reads and writes $n \times 8$ bytes = 32~MB, for a total of
$22 \times 64$~MB = 1.4~GB of device memory traffic per transform, leaving
little room for further compute scaling.

\subsection{Roofline Analysis and Occupancy}

To understand the compute vs.\ memory-bandwidth regime for each kernel, we apply
a simplified Roofline model to the T4.
The T4 delivers a peak FP32 throughput of approximately 8.1~TFLOPS and a peak
global memory bandwidth of 320~GB/s, giving an arithmetic intensity ridge point
of $\approx 25$~FLOP/byte. Our kernels operate on 64-bit integers rather than
floating-point, but the model is instructive as an order-of-magnitude guide.

The \texttt{modmul} kernel has near-perfect arithmetic intensity: each thread
reads two 8-byte operands and performs $\sim$5 integer operations (two 64-bit
multiplies, a shift, and a subtract), yielding an intensity of roughly $5/16
\approx 0.3$~ops/byte---well below the ridge point. However, because all
$n = 4$M threads are launched simultaneously and there is no global memory
traffic beyond the initial load of operands (each operand is used exactly once),
the kernel saturates compute rather than memory bandwidth: it is effectively
compute-bound at the warp level, with 2560 CUDA cores each handling independent
\texttt{modmul} chains with no shared-memory or inter-warp communication.

The NTT kernel regime shifts with $n$. At small $n$ ($\leq 2^{14}$), the working
set ($n \times 8$~bytes $\leq 128$~KB) fits in L1/shared memory across the 20
SMs, and the kernel is compute-bound. At $n = 2^{18}$ (2~MB), the working set
fits in T4's 4~MB L2 cache, maximising data reuse across the $\log_2 n = 18$
stage passes---this is where the 52$\times$ peak speedup is achieved. Beyond
$n = 2^{20}$ (8~MB), the working set overflows L2 and each stage requires a
full round-trip to global DRAM, making the kernel strongly memory-bandwidth-bound.
This explains why speedup drops from $52\times$ to $35\text{--}40\times$ at
$n \geq 2^{20}$ despite the additional parallelism.

Occupancy on sm\_75 is also a limiting factor at small $n$.
With 256 threads per block and $n/2$ total threads, the T4's 20~SMs are not
fully loaded below $n = 2^{13}$ (where the total thread count equals
$4096$ total threads $\ll$ $20 \times 2048 = 40960$ maximum resident threads). This partially explains the
sub-1$\times$ speedup at $n = 2^{10}$ independently of launch overhead.

\subsection{Barrett vs.\ Montgomery for NTT}

Montgomery multiplication can reduce the per-reduction operation count in
tight modexp loops, but for NTT, Barrett is the better choice: data stays in
the natural residue domain across all butterfly stages. Montgomery would require
converting the entire $n$-element array into Montgomery form before the NTT and
back out afterward---two additional O($n$) passes over device memory. With large
$n$ where the NTT is already memory-bandwidth-bound, these extra passes would
materially reduce effective speedup. Our engine is prime-agnostic: any $p$
requires only recomputing $\mu = \lfloor 2^k / p \rfloor$.

\subsection{Performance vs.\ Security Trade-offs}

Our implementation is optimized for throughput benchmarking and carries two
key security limitations.

\textbf{Timing side channel in modexp.}
The square-and-multiply implementation branches on exponent bits
(\texttt{if (e \& 1)}). If the exponent is secret (e.g., an RSA private key),
this creates a data-dependent timing channel. A production-safe deployment would
use a Montgomery ladder or always-square-and-multiply pattern with a constant
exponent.

\textbf{Data-dependent Barrett correction.}
The conditional subtraction \texttt{if (r >= p) r -= p} leaks one bit of the
residual. The branch-free equivalent
\texttt{r -= p \& -(u64)(r >= p)} eliminates this at negligible
throughput cost.

Beyond these two points, the GPU batch-parallelism model itself has
security relevance: the 277$\times$ speedup for batch modexp (1M independent
operations) reflects an ideal attacker-side advantage for brute-force or
parallel exhaustive search. A single large modular exponentiation (e.g., one
RSA-2048 decryption) remains inherently sequential; GPU advantage requires
batching. Table~\ref{tab:security} summarises the key security properties and
the recommended fixes for a production deployment.

\begin{table}[t]
\caption{Security Properties of the Implementation}
\label{tab:security}
\begin{center}
\begin{tabular}{|p{2.2cm}|p{2.0cm}|p{2.8cm}|}
\hline
\textbf{Property} & \textbf{Current status} & \textbf{Production fix} \\
\hline
modexp timing & Data-dependent (branching on exponent bits) & Montgomery ladder / always-square \\
\hline
Barrett correction & Data-dependent (\texttt{if r $\geq$ p}) & Branch-free \texttt{r -= p \& -(r$\geq$p)} \\
\hline
Field size & 30-bit (P32), 64-bit (P64) & 256-bit ECC or Kyber modulus ($q=3329$) \\
\hline
Batch parallelism & High throughput for both parties & Secret keys must not be batched in non-constant-time kernels \\
\hline
\end{tabular}
\end{center}
\end{table}

The field-size vs.\ throughput trade-off is measurable: moving from P32 (30-bit)
to P64 (64-bit) reduces modmul-dominated speedup from 569$\times$ to 277$\times$;
RSA-2048 or 256-bit ECC would require multi-limb arithmetic and further reduce
raw throughput in exchange for security margin.

%%--------------------------------------------------------------------
\section{Conclusion}
%%--------------------------------------------------------------------

We built a prime-agnostic GPU Barrett reduction engine in pure CUDA, validated
by a three-layer correctness protocol (GPU $\equiv$ CPU C++ $\equiv$
Python/sympy). Applied to the NTT, Barrett reduction achieves $52\times$ peak
speedup at the compute-memory balance point ($n=2^{18}$), while embarrassingly
parallel modmul reaches $569\times$ and batch Goldilocks modexp reaches
$277\times$. All results are obtained with zero correctness failures across
more than one million verified output elements.

Three engineering insights emerged from this project. \emph{First}, the
\texttt{\_\_host\_\_ \_\_device\_\_} dual-compilation strategy---validating 2M+
arithmetic cases locally on a CPU before touching a GPU---is the single most
effective tool for reducing debugging cost in GPU cryptography. \emph{Second},
Barrett reduction is strictly preferable to Montgomery for NTT workloads, because
it avoids two O($n$) domain-conversion memory passes between butterfly stages;
at large $n$ where the NTT is memory-bandwidth-bound, those passes would
materially reduce effective speedup. \emph{Third}, GPU speedup is not a fixed
property of a workload---it is a function of problem size, memory access pattern,
and kernel launch overhead. Our results show the GPU is actually \emph{slower}
than a single CPU core at $n=2^{10}$, underscoring that GPU acceleration requires
careful problem sizing.

The implementation is available in full at
\url{https://github.com/achirra-cipher/ece268-gpu-modarith}, including all source
files, the Colab benchmark notebook, and the Python oracle.

Future work includes: a negacyclic NTT variant for a direct Kyber/Dilithium
polynomial multiplication tie-in; a shared-memory tiled NTT kernel to eliminate
the per-stage global memory round-trip for small $n$; an OpenMP multi-threaded
CPU baseline for a fairer throughput comparison; and branch-free, constant-time
Barrett correction (\texttt{r -= p \& -(u64)(r >= p)}) for security-critical
deployment where timing side-channels are a concern.

\section*{Acknowledgment}

The authors thank the ECE 268 course staff at UC San Diego for project guidance
and Google Colab for providing T4 GPU access used in all benchmarking.

\begin{thebibliography}{00}

\bibitem{kyber}
R.~Avanzi et al., ``CRYSTALS-Kyber (version 3.02),''
NIST Post-Quantum Cryptography Finalist Specification, 2021.
\url{https://pq-crystals.org/kyber/}

\bibitem{dilithium}
L.~Ducas et al., ``CRYSTALS-Dilithium (version 3.1),''
NIST Post-Quantum Cryptography Finalist Specification, 2021.
\url{https://pq-crystals.org/dilithium/}

\bibitem{plonky2}
Polygon Zero Team, ``Plonky2: Fast Recursive Arguments with PLONK and FRI,''
2022. \url{https://github.com/mir-protocol/plonky2}

\bibitem{barrett}
P.~D. Barrett, ``Implementing the Rivest Shamir and Adleman public key
encryption algorithm on a standard digital signal processor,''
in \textit{Proc. CRYPTO}, Lecture Notes in Computer Science, vol.~263,
pp.~311--323, Springer, 1987.

\bibitem{montgomery}
P.~L. Montgomery, ``Modular multiplication without trial division,''
\textit{Mathematics of Computation}, vol.~44, no.~170, pp.~519--521, 1985.

\bibitem{nttref}
T.~P\"{o}ppelmann and T.~G\"{u}neysu, ``Towards practical lattice-based
public-key encryption on reconfigurable hardware,''
in \textit{Proc. SAC}, Lecture Notes in Computer Science, vol.~8282,
pp.~68--85, Springer, 2013.

\end{thebibliography}

\end{document}
